% Abstract: written last in practice; this is the working draft built from
% the four ideas in the author's Spanish outline (origin, what was done,
% the two regimes, what the tables leave unsaid). Target 250 words.
\begin{abstract}
Infrared correlation tables give, for each functional group, a wavenumber window and an adjective. They do not say how often a band is found in that window when the group is absent, and that is the number that decides whether the observation is evidence. We treated every published rule as a diagnostic test. A set of 101 rules for 60 functional groups, transcribed from one widely used interpretation course and frozen before any confirmatory analysis, was scored on 1,539 ATR spectra of synthesised compounds from the Chemotion repository and replicated on 5,124 transmission spectra from the NIST Chemistry WebBook, with the deposited structure as ground truth and sensitivity, specificity and likelihood ratios as the metrics. The main result is a division of the spectrum into two regimes. Above 1500~\wn{} most rules carry real evidence (median positive likelihood ratio 2.8 and 3.0 in the two collections, up to 43 for the saturated ester carbonyl); in the fingerprint region, at or below 1350~\wn{}, band position alone is close to uninformative (median 1.1 and 1.2), because almost every spectrum has a band in almost every window. The rankings agree across the two collections (Spearman 0.79; 87\% of rules within a factor of two). The audit also measures what the tables leave unsaid: intensity qualifiers cost more sensitivity than they buy specificity (the nitrile band reaches the stated intensity in one nitrile in ten), the benzene substitution patterns do not reach a likelihood ratio of 5 in either collection, and several rules are far more useful for ruling a group out than for ruling it in. We provide the quantified table, the frozen rule set and the scripts so that the audit can be extended to further rules and collections.
\end{abstract}

\noindent\textbf{Keywords:} infrared spectroscopy; correlation tables; group frequencies; functional group identification; likelihood ratio; diagnostic accuracy; open spectral data
